%%%%%%%%%%%%%%%%%%%%%%%%%%%%%%%%%%%%%%%%%%%%%%%%%%%%%%%%%%%%%%%%%%%%%%%%%%
%
%   structure.tex — оформление книги GOS-Book
%
%   Содержание файла:
%     1.  Геометрия страницы
%     2.  Базовые пакеты
%     3.  Шрифты и кодировки
%     4.  Цвета
%     5.  Предметный указатель
%     6.  Оглавление (основное и мини-ToC для частей)
%     7.  Колонтитулы (хедер и футер)
%     8.  Стили теорем и цветные рамки (defn, thm, lemm, ...)
%     9.  Замечания (remark)
%     10. Заголовки секций
%     11. Заголовки частей (\part)
%     12. Заголовки глав (\chapter)
%
%%%%%%%%%%%%%%%%%%%%%%%%%%%%%%%%%%%%%%%%%%%%%%%%%%%%%%%%%%%%%%%%%%%%%%%%%%


%========================================================================
% 1. ГЕОМЕТРИЯ СТРАНИЦЫ
%========================================================================

\usepackage[
  margin      = 1.6cm,
  top         = 1.5cm,
  bottom      = 1.8cm,
  left        = 1.9cm,
  right       = 1.9cm,
  headheight  = 31pt,
  headsep     = 0.4cm,
  footskip    = 1.2cm
]{geometry}


%========================================================================
% 2. БАЗОВЫЕ ПАКЕТЫ
%========================================================================

\usepackage{graphicx}                 % вставка картинок
\graphicspath{{pictures/}}            % каталог с картинками

\usepackage{etoolbox}                 % \patchcmd, \AtBeginEnvironment и т.п.
\usepackage{calc}                     % арифметика длин: \widthof, \heightof
\usepackage{lipsum}                   % заглушки текста (для тестов)

\usepackage{tikz}                     % рисование
\usetikzlibrary{calc, backgrounds}

\usepackage{enumitem}                 % кастомизация списков
\usepackage{booktabs}                 % красивые линейки в таблицах


%========================================================================
% 3. ШРИФТЫ И КОДИРОВКИ
%========================================================================

\usepackage{microtype}                % улучшение межбуквенных расстояний
\usepackage{cmap}                     % корректный поиск текста в PDF
\usepackage[utf8]{inputenc}           % входная кодировка
\usepackage[T2A]{fontenc}             % шрифтовая кодировка (кириллица)
\usepackage[english, russian]{babel}  % языки


%========================================================================
% 4. ЦВЕТА
%========================================================================

\usepackage{xcolor}

\definecolor{prpl}   {RGB}{150, 120, 182}  % фиолетовый — основной акцент
\definecolor{grn}    {RGB}{ 76, 153,  76}  % зелёный    — определения
\definecolor{chapclr}{RGB}{200,  60,  60}  % красный    — заголовки глав/билетов


%========================================================================
% 5. ПРЕДМЕТНЫЙ УКАЗАТЕЛЬ
%========================================================================
%
% Собирается через xindy. Команда для ручной сборки:
%     texindy -M mystyle.xdy -L russian -C utf8 GOS-Book.idx
% При использовании latexmk с .latexmkrc это делается автоматически.
%========================================================================

\usepackage[xindy]{imakeidx}

\indexsetup{level=\chapter}
\makeindex[program=texindy, options=-M mystyle.xdy -L russian -C utf8]

% \rindex — обёртка для \index, защищающая специальные символы (детокенизация).
\makeatletter
\newcommand{\rindex}[2][\imki@jobname]{%
  \index[#1]{\detokenize{#2}}%
}
\makeatother


%========================================================================
% 6. ОГЛАВЛЕНИЕ
%========================================================================

\usepackage{titletoc}

\contentsmargin{0cm}

% В основном ToC показываем только главы (без секций).
\addtocontents{toc}{\protect\setcounter{tocdepth}{0}}

% ---- Основное оглавление -------------------------------------------------

% Части
\titlecontents{part}[0cm]
  {\addvspace{5pt}\centering\large\bfseries}{}{}{}

% Главы
\titlecontents{chapter}[0.6cm]
  {\addvspace{5pt}\large\sffamily}
  {\contentslabel[\large\thecontentslabel]{0.6cm}}
  {}
  {\large\sffamily\titlerule*[.5pc]{.}\;\thecontentspage}

% Разделы
\titlecontents{section}[0.6cm]
  {\addvspace{3pt}\sffamily}
  {\contentslabel[\thecontentslabel]{0.6cm}}
  {}
  {\sffamily\titlerule*[.5pc]{.}\;\thecontentspage}
  []

% Подразделы
\titlecontents{subsection}[0.6cm]
  {\addvspace{1pt}\sffamily\small}
  {\contentslabel[\thecontentslabel]{1.25cm}}
  {}
  {\titlerule*[.5pc]{.}\;\thecontentspage}
  []

% Список рисунков
\titlecontents{figure}[0em]
  {\addvspace{-5pt}\sffamily}
  {\thecontentslabel\hspace*{1em}}
  {}
  {\titlerule*[.5pc]{.}\;\thecontentspage}
  []

% Список таблиц
\titlecontents{table}[0em]
  {\addvspace{-5pt}\sffamily}
  {\thecontentslabel\hspace*{1em}}
  {}
  {\titlerule*[.5pc]{.}\;\thecontentspage}
  []

% ---- Мини-оглавление на страницах частей --------------------------------

\addtocontents{ptc}{\protect\setcounter{tocdepth}{2}}

\titlecontents{lchapter}[0em]
  {\addvspace{10pt}\large\sffamily\bfseries}
  {\color{prpl}\contentslabel[\Large\thecontentslabel]{1.25cm}\color{prpl}}
  {}
  {\color{prpl}\large\sffamily\bfseries\;\titlerule*[.5pc]{.}\;\thecontentspage}

\titlecontents{lsection}[0em]
  {\sffamily\small}
  {\contentslabel[\thecontentslabel]{1.25cm}}
  {}
  {\sffamily\small\;\titlerule*[.5pc]{.}\;\thecontentspage}

\titlecontents{lsubsection}[.5em]
  {\normalfont\footnotesize\sffamily}
  {}
  {}
  {\normalfont\footnotesize\sffamily\;\titlerule*[.5pc]{.}\;\thecontentspage}


%========================================================================
% 7. КОЛОНТИТУЛЫ (ХЕДЕР И ФУТЕР)
%========================================================================
%
% Схема:
%   - Хедер:  название главы (чётные), название раздела (нечётные).
%             Без номеров страниц — чтобы не мешать длинным названиям.
%   - Футер:  номер страницы по внешнему краю.
%   - Всё в приглушённом фиолетовом цвете.
%========================================================================

\usepackage{fancyhdr}

\pagestyle{fancy}

% ---- Определение флага для appendix --------------------------------------

\makeatletter
\newcommand{\inappendix}{TT\fi\expandafter\ifx\@chapapp\appendixname}
\makeatother

% ---- Форматирование меток глав/секций ------------------------------------

\makeatletter
\renewcommand{\chaptermark}[1]{%
  \if\inappendix
    \markboth{#1}{}%
  \else
    \if@mainmatter
      \markboth{%
        \ifnum\value{secnumdepth}>-1
          \chaptertitlename\thechapter.\ #1%
        \else
          #1%
        \fi
      }{}%
    \else
      \markboth{#1}{}%
    \fi
  \fi
}
\makeatother

\renewcommand{\sectionmark}[1]{\markright{\thesection\hspace{5pt}#1}{}}

% ---- Наполнение хедера и футера -----------------------------------------

\fancyhf{}  % очищаем всё

\fancyhead[LO]{\nouppercase{\footnotesize\sffamily\color{prpl!80}\rightmark}}
\fancyhead[RE]{\nouppercase{\footnotesize\sffamily\color{prpl!80}\leftmark}}

\fancyfoot[LE]{\sffamily\small\color{prpl!80}\thepage}
\fancyfoot[RO]{\sffamily\small\color{prpl!80}\thepage}

\renewcommand{\headrulewidth}{0.4pt}
\renewcommand{\footrulewidth}{0pt}

% ---- Стиль plain (первые страницы глав) ---------------------------------

\fancypagestyle{plain}{%
  \fancyhf{}%
  \fancyfoot[LE]{\sffamily\small\color{prpl!80}\thepage}%
  \fancyfoot[RO]{\sffamily\small\color{prpl!80}\thepage}%
  \renewcommand{\headrulewidth}{0pt}%
  \renewcommand{\footrulewidth}{0pt}%
}

% ---- Цвет линии под хедером ---------------------------------------------

\makeatletter
\renewcommand{\headrule}{%
  {\color{prpl!40}\hrule\@height\headrulewidth\@width\headwidth}%
  \vskip-\headrulewidth
}
\makeatother

% Страницы предметного указателя тоже с fancy-стилем.
\indexsetup{othercode={\pagestyle{fancy}}}


%========================================================================
% 8. СТИЛИ ТЕОРЕМ И ЦВЕТНЫЕ РАМКИ
%========================================================================

\usepackage{amsmath, amsfonts, amssymb, amsthm}
\RequirePackage[framemethod=default]{mdframed}

% ---- Вспомогательные математические команды -----------------------------

\newcommand{\intoo}[2]{\mathopen{]}#1\,;#2\mathclose{[}}
\newcommand{\intff}[2]{\mathopen{[}#1\,;#2\mathclose{]}}
\newcommand{\ud}{\mathop{\mathrm{{}d}}\mathopen{}}

% ---- Стили теорем (amsthm) ----------------------------------------------

\makeatletter

% Заголовок фиолетовый, в рамке (для thm, lemm, stt, axiome, exerc)
\newtheoremstyle{prplnumbox}
  {0pt}{0pt}{\normalfont}{}
  {\small\bfseries\sffamily\color{prpl}}{.}{0.25em}
  {\small\sffamily\color{prpl}%
   \thmname{#1}\thmnumber{\@ifnotempty{#1}{}\@upn{\nobreakspace#2}}%
   \thmnote{\nobreakspace\sffamily\bfseries\color{black}(#3)}}

% Заголовок чёрный, в рамке (для defn, cons)
\newtheoremstyle{blacknumbox}
  {0pt}{0pt}{\normalfont}{}
  {\small\bfseries\sffamily}{.}{0.25em}
  {\small\sffamily%
   \thmname{#1}\nobreakspace\thmnumber{\@ifnotempty{#1}{}\@upn{#2}}%
   \thmnote{\nobreakspace\sffamily\bfseries(#3)}}

% Заголовок чёрный с квадратом-QED, без рамки (для exmpl)
\newtheoremstyle{blacknumex}
  {5pt}{5pt}{\normalfont}{}
  {\small\bfseries\sffamily}{.}{0.25em}
  {\small\sffamily{\qedsymbol}\nobreakspace%
   \thmname{#1}\nobreakspace\thmnumber{\@ifnotempty{#1}{}\@upn{#2}}%
   \thmnote{\nobreakspace\sffamily\bfseries(#3)}}

% Заголовок фиолетовый, без рамки
\newtheoremstyle{prplnum}
  {5pt}{5pt}{\normalfont}{}
  {\small\bfseries\sffamily\color{prpl}}{.}{0.25em}
  {\small\sffamily\color{prpl}%
   \thmname{#1}\nobreakspace\thmnumber{\@ifnotempty{#1}{}\@upn{#2}}%
   \thmnote{\nobreakspace\sffamily\bfseries(#3)}}

% Стиль для замечаний (простой, без нумерации)
\newtheoremstyle{note}
  {3pt}{3pt}{}{}
  {\bfseries}{.}{.5em}{}

\makeatother

\renewcommand{\qedsymbol}{$\blacksquare$}

% ---- Цветные рамки (mdframed) -------------------------------------------

% Замечание — плотная рамка со всех сторон.
\newmdenv[
  skipabove         = 7pt,
  skipbelow         = 7pt,
  backgroundcolor   = black!5,
  linecolor         = prpl,
  linewidth         = 1pt,
  innerleftmargin   = 5pt,
  innerrightmargin  = 5pt,
  innertopmargin    = 5pt,
  innerbottommargin = 5pt,
  leftmargin        = 0cm,
  rightmargin       = 0cm,
  splittopskip      = 7pt,
  splitbottomskip   = 7pt
]{tBox}

% Упражнение / теорема / лемма — жирная фиолетовая линия слева, светлая заливка.
\newmdenv[
  skipabove         = 7pt,
  skipbelow         = 7pt,
  rightline         = false,
  topline           = false,
  bottomline        = false,
  leftline          = true,
  linecolor         = prpl,
  linewidth         = 4pt,
  backgroundcolor   = prpl!10,
  innerleftmargin   = 5pt,
  innerrightmargin  = 5pt,
  innertopmargin    = 5pt,
  innerbottommargin = 5pt,
  leftmargin        = 0cm,
  rightmargin       = 0cm,
  splittopskip      = 7pt,
  splitbottomskip   = 7pt
]{eBox}

% Определение — жирная зелёная линия слева, без заливки.
\newmdenv[
  skipabove         = 7pt,
  skipbelow         = 7pt,
  rightline         = false,
  topline           = false,
  bottomline        = false,
  leftline          = true,
  linecolor         = grn,
  linewidth         = 4pt,
  innerleftmargin   = 5pt,
  innerrightmargin  = 5pt,
  innertopmargin    = 0pt,
  innerbottommargin = 0pt,
  leftmargin        = 0cm,
  rightmargin       = 0cm,
  splittopskip      = 7pt,
  splitbottomskip   = 7pt
]{dBox}

% Следствие — серая линия слева, светлая заливка.
\newmdenv[
  skipabove         = 7pt,
  skipbelow         = 7pt,
  rightline         = false,
  topline           = false,
  bottomline        = false,
  leftline          = true,
  linecolor         = gray,
  linewidth         = 4pt,
  backgroundcolor   = black!5,
  innerleftmargin   = 5pt,
  innerrightmargin  = 5pt,
  innertopmargin    = 5pt,
  innerbottommargin = 5pt,
  leftmargin        = 0cm,
  rightmargin       = 0cm,
  splittopskip      = 7pt,
  splitbottomskip   = 7pt
]{cBox}

% ---- Окружения теорем ---------------------------------------------------
%
% Для каждого типа определяются:
%   *T   — сам theorem-счётчик и стиль,
%   без T — окружение с обёрткой в цветную рамку.
% Штрихованные версии (defnn, thmn, ...) используют номер основного
% счётчика со штрихом: 3'
%
% Дополнительные обозначения:
%   \notation  — глобальный счётчик обозначений.
%   \notion    — замечание (без нумерации).
%========================================================================

\newtheorem{notation}{Обозначение}[chapter]

% Определение
\theoremstyle{blacknumbox}
\newtheorem{defnT}{Определение}[chapter]
\renewcommand{\thedefnT}{\arabic{defnT}}
\renewcommand{\theHdefnT}{\thechapter.\arabic{defnT}}
\newenvironment{defn}{\begin{dBox}\begin{defnT}}{\end{defnT}\end{dBox}}

\newtheorem{defnnT}{Определение}
\renewcommand{\thedefnnT}{\arabic{defnT}'}
\renewcommand{\theHdefnnT}{\thechapter.\arabic{defnnT}'}
\newenvironment{defnn}{\begin{dBox}\begin{defnnT}}{\end{defnnT}\end{dBox}}

% Теорема
\theoremstyle{prplnumbox}
\newtheorem{thmT}{Теорема}[chapter]
\renewcommand{\thethmT}{\arabic{thmT}}
\renewcommand{\theHthmT}{\thechapter.\arabic{thmT}}
\newenvironment{thm}{\begin{eBox}\begin{thmT}}{\end{thmT}\end{eBox}}

\newtheorem{thmnT}{Теорема}[chapter]
\renewcommand{\thethmnT}{\arabic{thmT}'}
\renewcommand{\theHthmnT}{\thechapter.\arabic{thmnT}'}
\newenvironment{thmn}{\begin{eBox}\begin{thmnT}}{\end{thmnT}\end{eBox}}

% Лемма
\newtheorem{lemmT}{Лемма}[chapter]
\renewcommand{\thelemmT}{\arabic{lemmT}}
\renewcommand{\theHlemmT}{\thechapter.\arabic{lemmT}}
\newenvironment{lemm}{\begin{eBox}\begin{lemmT}}{\end{lemmT}\end{eBox}}

\newtheorem{lemmnT}{Лемма}[chapter]
\renewcommand{\thelemmnT}{\arabic{lemmT}'}
\renewcommand{\theHlemmnT}{\thechapter.\arabic{lemmnT}'}
\newenvironment{lemmn}{\begin{eBox}\begin{lemmnT}}{\end{lemmnT}\end{eBox}}

% Утверждение
\newtheorem{sttT}{Утверждение}[chapter]
\renewcommand{\thesttT}{\arabic{sttT}}
\renewcommand{\theHsttT}{\thechapter.\arabic{sttT}}
\newenvironment{stt}{\begin{eBox}\begin{sttT}}{\end{sttT}\end{eBox}}

% Аксиома
\newtheorem{axiomeT}{Аксиома}[chapter]
\renewcommand{\theaxiomeT}{\arabic{axiomeT}}
\renewcommand{\theHaxiomeT}{\thechapter.\arabic{axiomeT}}
\newenvironment{axiome}{\begin{eBox}\begin{axiomeT}}{\end{axiomeT}\end{eBox}}

\newtheorem{axiomenT}{Аксиома}[chapter]
\renewcommand{\theaxiomenT}{\arabic{axiomeT}'}
\renewcommand{\theHaxiomenT}{\thechapter.\arabic{axiomenT}'}
\newenvironment{axiomen}{\begin{eBox}\begin{axiomenT}}{\end{axiomenT}\end{eBox}}

% Следствие
\theoremstyle{blacknumbox}
\newtheorem{consT}{Следствие}[chapter]
\renewcommand{\theconsT}{\arabic{consT}}
\renewcommand{\theHconsT}{\thechapter.\arabic{consT}}
\newenvironment{cons}{\begin{cBox}\begin{consT}}{\end{consT}\end{cBox}}

\newtheorem{consnT}{Следствие}[chapter]
\renewcommand{\theconsnT}{\arabic{consT}'}
\renewcommand{\theHconsnT}{\thechapter.\arabic{consnT}'}
\newenvironment{consn}{\begin{cBox}\begin{consnT}}{\end{consnT}\end{cBox}}

% Пример
\theoremstyle{blacknumex}
\newtheorem{exmplT}{Пример}[chapter]
\renewcommand{\theexmplT}{\arabic{exmplT}}
\renewcommand{\theHexmplT}{\thechapter.\arabic{exmplT}}
\newenvironment{exmpl}{\begin{exmplT}}{\end{exmplT}}

% Упражнение
\theoremstyle{prplnumbox}
\newtheorem{exercT}{Упражнение}[chapter]
\renewcommand{\theexercT}{\arabic{exercT}}
\renewcommand{\theHexercT}{\thechapter.\arabic{exercT}}
\newenvironment{exerc}{\begin{eBox}\begin{exercT}}{\end{exercT}\end{eBox}}

% Замечание (без номера)
\newtheorem*{notionT}{Замечание}
\newenvironment{notion}{\begin{tBox}\begin{notionT}}{\end{notionT}\end{tBox}}

% ---- Нумерация формул: по главам, полная (X.Y) --------------------------

\numberwithin{equation}{chapter}
\renewcommand{\theequation}{\thechapter.\arabic{equation}}
\renewcommand{\theHequation}{\thechapter.\arabic{equation}}

% ---- Решение (заголовок вместо стандартного "Proof") --------------------

\newenvironment{solution}
  {\begin{proof}[Решение.]}
  {\end{proof}}


%========================================================================
% 9. ЗАМЕЧАНИЕ (REMARK)
%========================================================================
%
% Абзац с отступами и кружком-меткой "R" слева. Используется как альтернатива
% notion, когда нужно выделить примечание без рамки.
%========================================================================

\newenvironment{remark}{%
  \par\vspace{10pt}\small
  \begin{list}{}{\leftmargin=35pt \rightmargin=25pt}%
  \item\ignorespaces
  \makebox[-2.5pt]{%
    \begin{tikzpicture}[overlay]
      \node[
        draw       = prpl!60,
        line width = 1pt,
        circle,
        fill       = prpl!25,
        font       = \sffamily\bfseries,
        inner sep  = 2pt,
        outer sep  = 0pt
      ] at (-15pt, 0pt) {\textcolor{prpl}{R}};
    \end{tikzpicture}%
  }%
  \advance\baselineskip -1pt
}{%
  \end{list}\vskip5pt
}


%========================================================================
% 10. ЗАГОЛОВКИ СЕКЦИЙ
%========================================================================
%
% Номер секции выносится на левое поле фиолетовым цветом.
%========================================================================

\makeatletter

\renewcommand{\@seccntformat}[1]{%
  \llap{\textcolor{prpl}{\csname the#1\endcsname}\hspace{1em}}%
}

\renewcommand{\section}{\@startsection{section}{1}{\z@}%
  {-4ex \@plus -1ex \@minus -0.4ex}%
  { 1ex \@plus  0.2ex}%
  {\normalfont\large\sffamily\bfseries}}

\renewcommand{\subsection}{\@startsection{subsection}{2}{\z@}%
  {-3ex \@plus -0.1ex \@minus -0.4ex}%
  {0.5ex \@plus  0.2ex}%
  {\normalfont\sffamily\bfseries}}

\renewcommand{\subsubsection}{\@startsection{subsubsection}{3}{\z@}%
  {-2ex \@plus -0.1ex \@minus -0.2ex}%
  {0.2ex \@plus  0.2ex}%
  {\normalfont\small\sffamily\bfseries}}

\renewcommand{\paragraph}{\@startsection{paragraph}{4}{\z@}%
  {-2ex \@plus -0.2ex \@minus 0.2ex}%
  {0.1ex}%
  {\normalfont\small\sffamily\bfseries}}

\makeatother


%========================================================================
% 11. ЗАГОЛОВКИ ЧАСТЕЙ (\part)
%========================================================================
%
% Премиальный минималистичный стиль:
%   - фон + водяной знак + линии на ВСЕХ страницах части (включая ToC)
%   - декор рисуется через eso-pic (\AddToShipoutPictureBG), НЕ через
%     pagestyle/fancyhdr — это устраняет утечку фона на страницы главы,
%     идущей следом, поскольку не зависит от того, как \chapter внутри
%     себя работает с \cleardoublepage/\thispagestyle.
%   - декор корректно завершается — не переходит на страницы глав
%   - ToC не пересекает верхнюю и нижнюю декоративные линии
%
% ТРЕБУЕТСЯ: \usepackage{eso-pic} в разделе 2 преамбулы.
%========================================================================

\newlength\esp   \setlength\esp{4pt}
\newlength\ecart \setlength\ecart{1.2cm - \esp}

\newcommand{\thepartimage}{}
\newcommand{\partimage}[1]{\renewcommand{\thepartimage}{#1}}

\makeatletter

% ---- Формат части в главном оглавлении ----------------------------------

\newcommand{\@mypartnumtocformat}[2]{%
  \setlength\fboxsep{0pt}%
  \noindent
  \colorbox{prpl!20}{%
    \strut\parbox[c][.7cm]{\ecart}{%
      \color{prpl!70}\Large\sffamily\bfseries\centering #1%
    }%
  }%
  \hskip\esp
  \colorbox{prpl!40}{%
    \strut\parbox[c][.7cm]{\linewidth-\ecart-\esp}{%
      \Large\sffamily\centering #2%
    }%
  }%
}

\newcommand{\@myparttocformat}[1]{%
  \setlength\fboxsep{0pt}%
  \noindent
  \colorbox{prpl!40}{%
    \strut\parbox[c][.7cm]{\linewidth}{%
      \Large\sffamily\centering #1%
    }%
  }%
}

% ---- Стили мини-оглавления для страниц частей ---------------------------
%
% Ширина боксов под номер увеличена, а отступ после номера сделан
% заметным (было 0.2em — слишком мало и почти не читалось на глаз,
% особенно когда номер занимал почти всю ширину бокса).

\titlecontents{lchapter}[2em]
  {\vspace{5pt}\normalsize\sffamily\bfseries\color{black}}
  {\color{prpl}\contentslabel[\thecontentslabel]{2em}\hspace{0.6em}\color{black}}
  {}
  {\titlerule*[5pt]{.}\color{prpl!80}\normalfont\sffamily\thecontentspage}
  [\vspace{0pt}]

\titlecontents{lsection}[3.2em]
  {\vspace{2pt}\footnotesize\sffamily\color{black!70}}
  {\contentslabel[\thecontentslabel]{2.4em}\hspace{0.6em}}
  {}
  {\titlerule*[5pt]{.}\color{black!50}\footnotesize\thecontentspage}

\titlecontents{lsubsection}[4em]
  {\vspace{1pt}\normalfont\footnotesize\sffamily\color{black!60}}
  {}
  {}
  {\titlerule*[5pt]{.}\footnotesize\color{black!50}\thecontentspage}

% ---- Размеры страницы части ---------------------------------------------

\newlength\partleftpad
\setlength\partleftpad{2cm}

\newlength\partrightpad
\setlength\partrightpad{2cm}

\newlength\partwatermarkw
\setlength\partwatermarkw{7cm}

\newlength\partwatermarkh
\setlength\partwatermarkh{8cm}

% Отступы текстового блока относительно декоративных линий (линии рисуются
% на 2cm сверху / 1.5cm снизу от физического края страницы).
\newlength\partcontenttopskip
\setlength\partcontenttopskip{2.5cm}

\newlength\partcontentbottomskip
\setlength\partcontentbottomskip{2cm}

% Левое/правое поле — как в основной геометрии документа (раздел 1).
\newlength\partsidemargin
\setlength\partsidemargin{1.9cm}

% ---- Функции рисования декоративного фона --------------------------------
%
% current page — встроенный якорь TikZ, не требует межстраничной "памяти";
% remember picture тут формально не нужен, но безвреден — оставлен для
% совместимости, если где-то ещё в проекте есть перекрёстные ссылки.

\newcommand{\@drawpartbackground}{%
  \begin{tikzpicture}[remember picture, overlay]
    \fill[prpl!5]
      (current page.north west) rectangle (current page.south east);
    \node[
      anchor    = north east,
      inner sep = 0pt,
      outer sep = 0pt,
      opacity   = 0.10
    ] at ([xshift=-1.2cm, yshift=-2.5cm]current page.north east) {%
      \makebox[\partwatermarkw][c]{%
        \resizebox{!}{\partwatermarkh}{%
          \rmfamily\bfseries\color{prpl}\thepart%
        }%
      }%
    };
    \draw[prpl!50, line width=0.8pt]
      ([xshift=\partleftpad,  yshift=-2cm]current page.north west) --
      ([xshift=-\partrightpad, yshift=-2cm]current page.north east);
    \draw[prpl!50, line width=0.8pt]
      ([xshift=\partleftpad,  yshift=1.5cm]current page.south west) --
      ([xshift=-\partrightpad, yshift=1.5cm]current page.south east);
  \end{tikzpicture}%
}

\newcommand{\@drawspartbackground}{%
  \begin{tikzpicture}[remember picture, overlay]
    \fill[prpl!5]
      (current page.north west) rectangle (current page.south east);
    \draw[prpl!50, line width=0.8pt]
      ([xshift=\partleftpad,  yshift=-2cm]current page.north west) --
      ([xshift=-\partrightpad, yshift=-2cm]current page.north east);
    \draw[prpl!50, line width=0.8pt]
      ([xshift=\partleftpad,  yshift=1.5cm]current page.south west) --
      ([xshift=-\partrightpad, yshift=1.5cm]current page.south east);
  \end{tikzpicture}%
}

% ---- Вкл/выкл режима "страница части" ------------------------------------
%
% Декор регистрируется через \AddToShipoutPictureBG — он вешается
% на КАЖДЫЙ \shipout, пока не будет вызван \ClearShipoutPictureBG.
% Это устраняет зависимость от pagestyle/thispagestyle: даже если
% \chapter вставит дополнительную пустую страницу через свой
% \cleardoublepage, эта страница просто не будет декорирована, если
% мы уже вызвали \ClearShipoutPictureBG к этому моменту.
%
% pagestyle{empty} на страницах части — нет ни бегущих колонтитулов,
% ни номеров страниц (как и было в исходном варианте с partpage).
%========================================================================

\newcommand{\@partpageon}[1]{%
  \pagestyle{empty}%
  \thispagestyle{empty}%
  \AddToShipoutPictureBG{#1}%
}

\newcommand{\@partpageoff}{%
  \ClearShipoutPictureBG
  \pagestyle{fancy}%
}

\newcommand{\@partmodeon}{%
  \newgeometry{%
    top        = \partcontenttopskip,
    bottom     = \partcontentbottomskip,
    left       = \partsidemargin,
    right      = \partsidemargin,
    headheight = 0pt,
    headsep    = 0pt,
    footskip   = 0pt
  }%
}

\newcommand{\@partmodeoff}{%
  \restoregeometry
}

% ---- Переопределение \part -----------------------------------------------
%
% Стандартный \part из book.cls сам вставляет \null\vfil перед вызовом
% \@part/\@spart — эта пружина в паре с \vfil в конце \@endpart
% центрировала контент по вертикали вместо прижатия к верху ("съезжание"
% заголовка). Оставляем только переход на страницу + режим колонок.
%
% \cleardoublepage здесь — стандартное поведение книжной вёрстки:
% часть всегда открывается на нечётной (правой) странице; если предыдущая
% страница чётная, перед частью появится пустой лист-разворот — это
% нормально для двусторонней печати. Если вам это не нужно вовсе,
% замените \cleardoublepage на \clearpage ниже.
%========================================================================

\renewcommand\part{%
  \if@openright
    \cleardoublepage
  \else
    \clearpage
  \fi
  \if@twocolumn
    \onecolumn
    \@tempswatrue
  \else
    \@tempswafalse
  \fi
  \secdef\@part\@spart}

% ---- Нумерованная часть -------------------------------------------------

\def\@part[#1]#2{%
  \ifnum \c@secnumdepth >-2\relax
    \refstepcounter{part}%
    \addcontentsline{toc}{part}{%
      \texorpdfstring{%
        \protect\@mypartnumtocformat{\thepart}{#1}%
      }{%
        \partname~\thepart\ ---\ #1%
      }%
    }%
  \else
    \addcontentsline{toc}{part}{%
      \texorpdfstring{\protect\@myparttocformat{#1}}{#1}%
    }%
  \fi
  \startcontents
  \markboth{}{}%
  \@partpageon{\@drawpartbackground}%
  \@partmodeon
  \begingroup
    \parindent=0pt
    {\sffamily\bfseries\color{prpl!80}\normalsize
     \MakeUppercase{\partname}\ \thepart\par}
    \vspace{0.5cm}
    {\sffamily\bfseries\color{black}\Huge
     \raggedright #2\par}
    \vspace{1.3cm}
    {\color{prpl!50}\hrule height 0.4pt}
    \vspace{0.3cm}
    {\sffamily\bfseries\color{prpl!80}\small
     \MakeUppercase{Содержание части}\par}
    \vspace{0.3cm}
    \printcontents{l}{0}{\setcounter{tocdepth}{1}}%
  \endgroup
  \@endpart
}

% ---- Ненумерованная часть -----------------------------------------------

\def\@spart#1{%
  \startcontents
  \phantomsection
  \@partpageon{\@drawspartbackground}%
  \@partmodeon
  \begingroup
    \parindent=0pt
    {\sffamily\bfseries\color{prpl!80}\normalsize
     \MakeUppercase{\partname}\par}
    \vspace{0.5cm}
    {\sffamily\bfseries\color{black}\Huge
     \raggedright #1\par}
  \endgroup
  \addcontentsline{toc}{part}{%
    \texorpdfstring{%
      \setlength\fboxsep{0pt}%
      \noindent
      \protect\colorbox{prpl!40}{%
        \strut\protect\parbox[c][.7cm]{\linewidth}{%
          \Large\sffamily\protect\centering #1\quad\mbox{}%
        }%
      }%
    }{#1}%
  }%
  \@endpart
}

% ---- Завершение части ---------------------------------------------------
%
% \clearpage завершает все страницы части (первая + продолжения ToC) —
% на этот момент фон ЕЩЁ зарегистрирован через eso-pic, поэтому все они
% декорируются корректно. Сразу после — \ClearShipoutPictureBG снимает
% регистрацию фона и возвращает обычный pagestyle. Никакой ручной логики
% "открыть страницу вправо" здесь больше нет — это делает сам \chapter
% через свой \cleardoublepage; дублирование этой логики и было причиной
% лишних/пропадающих пустых листов в прошлых версиях.
%========================================================================

\def\@endpart{%
  \vfil
  \clearpage
  \@partmodeoff
  \@partpageoff
  \if@tempswa
    \twocolumn
  \fi}

\makeatother


%========================================================================
% 12. ЗАГОЛОВКИ ГЛАВ (\chapter)
%========================================================================
%
% Минималистичный современный стиль:
%   строка "Глава N" (фиолетовым)
%   крупное название главы (чёрным)
%   тонкая фиолетовая линия во всю ширину текста
%========================================================================

\makeatletter

% Нумерованные главы.
\def\@makechapterhead#1{%
  \vspace*{10pt}%
  {\parindent\z@ \raggedright \normalfont
   \if@mainmatter
     {\sffamily\bfseries\large\color{prpl}%
      \chaptertitlename\ \thechapter}%
     \par\vspace{4pt}%
     {\sffamily\bfseries\Huge\color{black}#1\par}%
   \else
     {\sffamily\bfseries\Huge\color{black}#1\par}%
   \fi
   \vspace{8pt}%
   {\color{prpl}\rule{\textwidth}{1.2pt}}%
   \par\vspace{20pt}%
  }%
}

% Ненумерованные главы (\chapter*).
\def\@makeschapterhead#1{%
  \vspace*{10pt}%
  {\parindent\z@ \raggedright \normalfont
   {\sffamily\bfseries\Huge\color{black}#1\par}%
   \vspace{8pt}%
   {\color{prpl}\rule{\textwidth}{1.2pt}}%
   \par\vspace{20pt}%
  }%
}

\makeatother